% The environments and macros of the modules (modules/*/*.qpp), for every manifest that prints them: the frame of
% a paper reads this file with \input{../../modules/macros} and keeps only its own title, abstract and bibliography
% style. Add a macro here, never in a frame: .github/scripts/check_printed.py refuses a control sequence of the printed
% modules that is neither defined here nor a command of LaTeX or of a package of the frame (its list STANDARD).
% the kinds that number with the theorems count on an alias of their counter (aliascnt), so that a reference names its
% own kind (Lemma, Remark, ...) and not the kind of the counter
\usepackage{aliascnt}
\theoremstyle{plain}
\newtheorem{theorem}{Theorem}[section]
\newaliascnt{lemma}{theorem}\newtheorem{lemma}[lemma]{Lemma}\aliascntresetthe{lemma}
\newaliascnt{proposition}{theorem}\newtheorem{proposition}[proposition]{Proposition}\aliascntresetthe{proposition}
\newaliascnt{corollary}{theorem}\newtheorem{corollary}[corollary]{Corollary}\aliascntresetthe{corollary}
\newaliascnt{conjecture}{theorem}\newtheorem{conjecture}[conjecture]{Conjecture}\aliascntresetthe{conjecture}
\newtheorem{fact}{Fact}
\renewcommand{\thefact}{F\arabic{fact}}
\theoremstyle{definition}
\newaliascnt{definition}{theorem}\newtheorem{definition}[definition]{Definition}\aliascntresetthe{definition}
\newaliascnt{example}{theorem}\newtheorem{example}[example]{Example}\aliascntresetthe{example}
\newaliascnt{remark}{theorem}\newtheorem{remark}[remark]{Remark}\aliascntresetthe{remark}
\newaliascnt{question}{theorem}\newtheorem{question}[question]{Question}\aliascntresetthe{question}
\newtheorem{problem}{Problem}
\newtheorem{construction}{Construction}
\crefname{fact}{Fact}{Facts}
\crefname{construction}{Construction}{Constructions}
\crefname{problem}{Problem}{Problems}
\crefname{conjecture}{Conjecture}{Conjectures}
\crefname{question}{Question}{Questions}

\newcommand{\Q}{\mathbb{Q}}
\newcommand{\Z}{\mathbb{Z}}
\newcommand{\N}{\mathbb{N}}
\newcommand{\R}{\mathbb{R}}
\newcommand{\C}{\mathbb{C}}
\newcommand{\Fp}{\mathbb F_p}
\newcommand{\Zp}{\Z_{(p)}}
\newcommand{\bw}{\mathbf{w}}
\newcommand{\bu}{\mathbf{u}}
\newcommand{\bv}{\mathbf{v}}
\newcommand{\bz}{\mathbf{z}}
\newcommand{\bx}{\mathbf{x}}
\newcommand{\by}{\mathbf{y}}
\newcommand{\bp}{\mathbf{p}}
\newcommand{\bq}{\mathbf{q}}
\newcommand{\bh}{\mathbf{h}}
\newcommand{\ba}{\mathbf{a}}
\newcommand{\bb}{\mathbf{b}}
\newcommand{\bd}{{\boldsymbol{\delta}}}
\newcommand{\bxi}{{\boldsymbol{\xi}}}
\newcommand{\balpha}{{\boldsymbol{\alpha}}}
\newcommand{\bbeta}{{\boldsymbol{\beta}}}
\newcommand{\one}{\mathbf{1}}
% the halves of the bideterminant: of a couple of tables (T, U) standing for T - U, and of one table (Gondran and
% Minoux; Akian, Gaubert and Guterman, arXiv:0812.3496, Def. 4.14 and Remark 4.17)
\newcommand{\bidetP}[2]{\lvert(#1,#2)\rvert^{+}}
\newcommand{\bidetM}[2]{\lvert(#1,#2)\rvert^{-}}
\newcommand{\bidetPa}[1]{\lvert #1\rvert^{+}}
\newcommand{\bidetMa}[1]{\lvert #1\rvert^{-}}
\newcommand{\Kb}{\bar K}
\newcommand{\Leb}{\operatorname{Leb}}
\DeclareMathOperator{\trop}{trop}
\DeclareMathOperator{\conv}{conv}
\DeclareMathOperator{\rec}{rec}
\DeclareMathOperator{\spn}{span}
\DeclareMathOperator{\supp}{supp}
\DeclareMathOperator{\adj}{adj}
\DeclareMathOperator{\sgn}{sgn}
\DeclareMathOperator{\ac}{ac}
\DeclareMathOperator{\val}{val}
\newcommand{\cont}{\operatorname{cont}}
\DeclareMathOperator{\re}{Re}
\DeclareMathOperator{\im}{Im}
\newcommand{\bt}{\mathbf{t}}
\newcommand{\bc}{\mathbf{c}}
\newcommand{\bs}{\mathbf{s}}
\newcommand{\bl}{\mathbf{l}}
\newcommand{\be}{{\boldsymbol{\varepsilon}}}
\newcommand{\btau}{{\boldsymbol{\tau}}}
\newcommand{\btheta}{{\boldsymbol{\theta}}}
\newcommand{\bsigma}{{\boldsymbol{\sigma}}}
\newcommand{\bpsi}{{\boldsymbol{\psi}}}
\newcommand{\bchi}{{\boldsymbol{\chi}}}
\newcommand{\vol}{\operatorname{vol}}
\DeclareMathOperator{\tr}{tr}
\DeclareMathOperator{\dist}{dist}
% the named expressions of the proof syntax (def:syntax): the generators, the patterns, the constants
\newcommand{\gI}{\mathrm{I}}
\newcommand{\gV}{\mathrm{V}}
\newcommand{\gEq}{\mathrm{eq}}
\newcommand{\gInd}{\mathrm{ind}}
% a fixture of a statement item (qpp/README.md, Fixtures); the article prints nothing for it
\newcommand{\FIX}[1]{}
\newcommand{\Pt}[1]{\mathrm{P#1}}
\newcommand{\Kt}{\mathrm{K}}
\newcommand{\Rf}[1]{\textup{(R#1)}}
\newcommand{\Mf}{\textup{(M)}}
% the syntax of a proof: after one or two sentences of orientation, the display expr, one line
% \eit{LABEL}{EXPRESSION} per item, condition or case, the expression in math mode; its commas and semicolons allow
% a line break; a line breaks at a semicolon or at a space \ (as after the comma of a list), never inside f(x,y)
{\catcode`\;=\active\gdef;{\mathchar"603B\penalty100\relax}}
\expandafter\let\expandafter\pcctrlspace\csname\space\endcsname
\def\pcbreakspace{\ifmmode\penalty200\fi\pcctrlspace}
\newenvironment{expr}{\par\smallskip\raggedright\mathcode`\;="8000\let\ \pcbreakspace\relax}{\qedhere\par}
\newcommand{\eit}[2]{\par\hangindent=1.5em\hangafter=1\noindent\textup{#1}\nobreak\quad\(#2\)}
% the syntax of a statement, a definition or a construction: after one sentence of description, the
% display stmt, one line \sit{LABEL}{CONTENT} per hypothesis, item, condition or object, the content in math mode
% (numbered equations may stand between the lines)
\newenvironment{stmt}{\par\smallskip\raggedright\mathcode`\;="8000\let\ \pcbreakspace\relax}{\par\smallskip}
\newcommand{\sit}[2]{\par\hangindent=1.5em\hangafter=1\noindent\textup{\textsc{#1}}\nobreak\quad\(#2\)}
% a row of a statement that is prose (the origin of a result), in text mode, so that it breaks across lines
\newcommand{\tit}[2]{\par\hangindent=1.5em\hangafter=1\noindent\textup{\textsc{#1}}\nobreak\quad #2}
% a proof transcoded from its words: one paragraph \pline{LABEL}{SENTENCES} per line of the proof, in text mode
\newenvironment{plines}{\par\smallskip\raggedright}{\qedhere\par}
\newcommand{\pline}[2]{\par\hangindent=1.5em\hangafter=1\noindent\textup{#1}\nobreak\quad #2}
